% SPDX-License-Identifier: CC-BY-4.0
\section{The lower bound at every origin order}\label{sec:lower}

Throughout this section \(\sigma\) is finite with \(|\sigma|=n\ge1\). Besides
\(\Ft[x_\sigma]\) we use a second polynomial ring \(\Ft[y_\sigma]\) in independent
variables \(y_i\). For \(A\in\Ft[y_\sigma]\), \(A(x^2+x)\) denotes the substitution
\(y_i\mapsto x_i^2+x_i\). For \(\varepsilon\subseteq\sigma\) put
\(x^\varepsilon=\prod_{i\in\varepsilon}x_i\), and for an exponent vector \(e\) put
\(|e|=\sum_ie_i\). Recall that \(f_d\) is the homogeneous component of degree
\(d\) of \(f\), and put \(f_d=0\) for \(d<0\). For \(K\ge0\) let
\[
 \hat x_K(Y)=\sum_{j<K}Y^{2^j}.
\]

\begin{lemma}[The local inverse]\label{lem:local-inverse}
Over \(\Ft\), \(\hat x_K(Y)^2+\hat x_K(Y)=Y+Y^{2^K}\). The series
\(\hat x(Y)=\sum_{j\ge0}Y^{2^j}\) satisfies \(\hat x^2+\hat x=Y\) in \(\Ft[[Y]]\), and
\(\hat x_K\) is its truncation below degree \(2^K\)
\textnormal{\small\leanref{claims/PerOrderYCoefficient.lean}{[Lean: \(\hat x_K\)]}\,\leanref{claims/ArtinSchreierSeries.lean}{[Lean: series]}}.
\end{lemma}

\begin{proof}
By Frobenius, \(\hat x_K^2+\hat x_K=\sum_{j<K}(Y^{2^{j+1}}+Y^{2^j})\), which
telescopes to \(Y+Y^{2^K}\). The
coefficients of \(\hat x\) and \(\hat x_K\) agree below degree \(2^K\), so the
coefficients of \(\hat x^2+\hat x\) and \(\hat x_K^2+\hat x_K\) agree below degree
\(2^K\); letting \(K\) grow gives \(\hat x^2+\hat x=Y\).
\end{proof}

So \(\hat x_K(y)\) inverts \(x\mapsto x^2+x\) up to order \(2^K\): it is the
truncation of the additive root of the Artin--Schreier equation \(X^2+X=Y\).

\paragraph{Plan.} Suppose \(\deg P<\Phi\). Lemmas~\ref{lem:expansion}
and~\ref{lem:congruence} turn the degree bound and the multiplicities into the
vanishing of certain homogeneous parts of \(w\prod_{i\in T}\hat x_K(y_i)\), and
Lemma~\ref{lem:lowest-part} shows that \(w\) starts in degree exactly \(\ell\).
These parts are values, at the generators \(y_i\), of additive forms
(Lemma~\ref{lem:forms}), and a recursion carries the vanishing from tuples of
distinct generators to all tuples, hence to the span \(V\). Vanishing on \(V\) is
a statement about values, not coefficients. The reduction modulo
\(\prod_{v\in V}(Y-v)\) (Lemma~\ref{lem:grid-reduction}) turns it into a
statement about coefficients, and one coefficient is \(w_\ell\).

\subsection{The expansion}

\begin{lemma}[Expansion]\label{lem:expansion}
Every \(P\in\Ft[x_\sigma]\) can be written as
\[
 P=\sum_{\varepsilon\subseteq\sigma}x^\varepsilon A_\varepsilon(x^2+x),\qquad
 A_\varepsilon\in\Ft[y_\sigma].
\]
For every such expansion, if the monomial \(y^e\) occurs in \(A_\varepsilon\), then
\(|\varepsilon|+2|e|\le\deg P\) \lean{claims/PerOrderExpansion.lean}.
\end{lemma}

\begin{proof}
For existence it suffices to expand a monomial, and by multiplicativity a power
\(x_i^b\). By induction on \(b\), \(x_i^b\) is a sum of terms \(x_i^\eta y_i^c\) with
\(\eta\in\{0,1\}\): this holds for \(b\le1\), and
\(x_i^b=x_i^{b-2}(x_i^2+x_i)+x_i^{b-1}\).

For the degree bound, let \(W\) be the largest value of \(|\varepsilon|+2|e|\) over
the pairs \((\varepsilon,e)\) with \(y^e\) in \(A_\varepsilon\), with coefficient
\(c_{\varepsilon,e}\). The polynomial \(x^\varepsilon(x^2+x)^e\) has degree
\(|\varepsilon|+2|e|\), and its homogeneous component of that degree is the single
monomial \(x^{\mathbf 1_\varepsilon+2e}\). The map \((\varepsilon,e)\mapsto\mathbf
1_\varepsilon+2e\) is injective, since \(\varepsilon\) is the set of odd
coordinates. So the component of degree \(W\) of \(P\) is
\(\sum c_{\varepsilon,e}x^{\mathbf 1_\varepsilon+2e}\) over the pairs of weight
\(W\), a sum of distinct monomials, which is nonzero. Hence \(\deg P\ge W\).
\end{proof}

\subsection{The multiplicity conditions}

\begin{lemma}[Congruence]\label{lem:congruence}
Let \(k\ge1\) and \(K\) with \(2^K\ge k\). Let \(P\in\Ft[x_\sigma]\) have
\(\mult_a(P)\ge k\) for every nonzero \(a\in\Ft^\sigma\), and let
\(P=\sum_\varepsilon x^\varepsilon A_\varepsilon(x^2+x)\) be any expansion. Put
\(\hat x_i=\hat x_K(y_i)\) and \(w=P(\hat x)\in\Ft[y_\sigma]\). Then for every
\(\varepsilon\subseteq\sigma\) and every \(d<k\),
\[
 (A_\varepsilon)_d=\Bigl(w\prod_{i\notin\varepsilon}(1+\hat x_i)\Bigr)_d
\]
\lean{claims/PerOrderYCoefficient.lean}.
\end{lemma}

\begin{proof}
Write \(\equiv\) for equality of all homogeneous components of degree below \(k\).
\begin{itemize}
 \item \emph{The substitution.} For \(\gamma\in\Ft^\sigma\) substitute
 \(x_i\mapsto\hat x_i+\gamma_i\). Since \(\gamma_i^2=\gamma_i\) and the cross term
 vanishes in characteristic two, \(x_i^2+x_i\) goes to
 \(\hat x_i^2+\hat x_i=y_i+y_i^{2^K}=:\hat y_i\) for both values of \(\gamma_i\)
 (Lemma~\ref{lem:local-inverse}). Hence
 \[
  P(\hat x+\gamma)=\sum_{\varepsilon'}\prod_{i\in\varepsilon'}(\hat x_i+\gamma_i)\,
  A_{\varepsilon'}(\hat y).
 \]
 \item \emph{Nonzero shifts.} The \(\hat x_i\) have no constant term, so for
 \(\gamma\ne\zero\) Lemma~\ref{lem:substitution} and \(\mult_\gamma(P)\ge k\) give
 \(\operatorname{ord}P(\gamma+\hat x)\ge k\), that is, \(P(\hat x+\gamma)\equiv0\).
 \item \emph{Inversion.} For \(c\in\Ft\) put \(N_i(c)=1\) if \(i\in\varepsilon\), and
 otherwise \(N_i(0)=1+\hat x_i\), \(N_i(1)=\hat x_i\). A check of the four cases
 gives
 \[
  \sum_{c\in\Ft}N_i(c)\,(\hat x_i+c)^{[i\in\varepsilon']}
  =\bigl[\,i\in\varepsilon\Leftrightarrow i\in\varepsilon'\,\bigr].
 \]
 For instance, for \(i\notin\varepsilon\) and \(i\in\varepsilon'\) the sum is
 \((1+\hat x_i)\hat x_i+\hat x_i(\hat x_i+1)=0\). Multiplying over \(i\) and summing
 over \(\gamma\in\Ft^\sigma\) gives
 \(\sum_\gamma\prod_iN_i(\gamma_i)\prod_{i\in\varepsilon'}(\hat x_i+\gamma_i)
 =[\varepsilon'=\varepsilon]\), and therefore
 \[
  A_\varepsilon(\hat y)=\sum_{\gamma}\prod_iN_i(\gamma_i)\,P(\hat x+\gamma).
 \]
 \item \emph{Conclusion.} Modulo \(\equiv\) only \(\gamma=\zero\) survives, and it
 contributes \(\prod_{i\notin\varepsilon}(1+\hat x_i)\,w\). Finally
 \(\hat y_i-y_i=y_i^{2^K}\) has order \(2^K\ge k\), so
 \(A_\varepsilon(\hat y)\equiv A_\varepsilon(y)\). \qedhere
\end{itemize}
\end{proof}

\begin{lemma}[The lowest part]\label{lem:lowest-part}
Let \(0\le\ell<k\) and \(P\in\Ft[x_\sigma]\) with \(\mult_{\zero}(P)=\ell\), and put
\(w=P(\hat x_k(y_1),\ldots,\hat x_k(y_n))\). Then \(w_d=0\) for \(d<\ell\) and
\(w_\ell\ne0\) \lean{claims/PerOrderLowerBound.lean}.
\end{lemma}

\noindent No condition at the nonzero points is needed here. In the proof,
\(K=k\), so \(2^K\ge k\).

\begin{proof}
Since the \(\hat x_i\) have no constant term, \(\operatorname{ord}w\ge\ell\)
(Lemma~\ref{lem:substitution}). Substitute \(y_i\mapsto x_i^2+x_i\) in \(w\). By
Lemma~\ref{lem:telescope} with \(t=x_i\), \(\hat x_K(x_i^2+x_i)=x_i+x_i^{2^K}\), so
the result is \(P(x+x^{2^K})\). The difference \(P(x+x^{2^K})-P(x)\) lies in the
ideal generated by the \(x_i^{2^K}\), so it has order at least \(2^K\ge k>\ell\),
and \(P(x+x^{2^K})\) has order exactly \(\ell\). If \(w_\ell\) were zero, then
\(\operatorname{ord}w>\ell\), and since each \(x_i^2+x_i\) has order one,
\(P(x+x^{2^K})\) would have order above \(\ell\), by Lemma~\ref{lem:substitution}.
\end{proof}

\subsection{Additive forms}

Fix \(k\ge1\) and \(w\in\Ft[y_\sigma]\). For \(t\ge0\), \(j\ge0\) and
\(Y_1,\ldots,Y_j\in\Ft[y_\sigma]\) put
\[
 E^{(t)}_j(Y_1,\ldots,Y_j)=\sum_{\tau\in\{0,\ldots,k-1\}^j}
 Y_1^{2^{\tau_1}}\cdots Y_j^{2^{\tau_j}}\;w_{k-1-t-\sum_a2^{\tau_a}} .
\]

\begin{lemma}[Forms]\label{lem:forms}
For all \(t,j\ge0\):
\begin{enumerate}
 \item \(E^{(t)}_j\) is symmetric, and additive in each argument.
 \item For \(t<k\) and \(i_1,\ldots,i_j\in\sigma\),
 \(E^{(t)}_j(y_{i_1},\ldots,y_{i_j})=\bigl(w\prod_a\hat x_k(y_{i_a})\bigr)_{k-1-t}\).
 \item For \(j\ge0\) and all \(Y,Z_1,\ldots,Z_j\),
 \[
  E^{(t)}_{j+2}(Y,Y,Z)=E^{(t)}_{j+1}(Y,Z)+Y\,E^{(t+1)}_j(Z).
 \]
\end{enumerate}
\lean{claims/PerOrderForms.lean}
\end{lemma}

\noindent Part~2 is stated for \(t<k\), the only case used; this also keeps the
degree \(k-1-t\) nonnegative in the Lean statement.

\begin{proof}
\begin{enumerate}
 \item Permuting the arguments permutes the index set. Additivity holds because
 \(Y\mapsto Y^{2^s}\) is additive in characteristic two.
 \item Expand \(\prod_a\hat x_k(y_{i_a})=\sum_\tau\prod_ay_{i_a}^{2^{\tau_a}}\) over
 \(\tau\in\{0,\ldots,k-1\}^j\). Each product is homogeneous of degree
 \(\sum_a2^{\tau_a}\), so the component of degree \(k-1-t\) of \(w\) times it is
 \(\prod_ay_{i_a}^{2^{\tau_a}}w_{k-1-t-\sum_a2^{\tau_a}}\).
 \item In \(E^{(t)}_{j+2}(Y,Y,Z)\) the terms with \(\tau_1\ne\tau_2\) cancel in pairs
 under the exchange of \(\tau_1\) and \(\tau_2\). The terms with \(\tau_1=\tau_2=\nu\)
 give \(Y^{2^{\nu+1}}w_{k-1-t-2^{\nu+1}-\Sigma}\), where \(\Sigma\) is the contribution of
 \(Z\). In \(E^{(t)}_{j+1}(Y,Z)\), the terms with \(Y\)-exponent \(2^{\nu+1}\),
 \(0\le\nu<k-1\), are the same terms. The remaining diagonal term \(\nu=k-1\) is zero,
 because \(2^k>k-1-t\) makes the index of \(w\) negative. The terms of
 \(E^{(t)}_{j+1}(Y,Z)\) with \(Y\)-exponent \(1\) are
 \(Y\,w_{k-1-(t+1)-\Sigma}\), which sum to \(Y\,E^{(t+1)}_j(Z)\). Signs do not
 matter in characteristic two. \qedhere
\end{enumerate}
\end{proof}

\subsection{Reduction on a grid}

\begin{lemma}[Grid vanishing {\cite[Lemma~2.1]{Alo99}}]\label{lem:grid-vanishing}
Let \(D\) be an integral domain, \(V\subseteq D\) finite, and
\(R\in D[Y_1,\ldots,Y_h]\) of degree less than \(|V|\) in each variable. If \(R\)
vanishes on \(V^h\), then \(R=0\).
\end{lemma}

\begin{proof}
Induct on \(h\). Write \(R=\sum_{e<|V|}R_e(Y_1,\ldots,Y_{h-1})Y_h^e\). For fixed
\(v'\in V^{h-1}\), the polynomial \(R(v',Y_h)\) has degree less than \(|V|\) and
\(|V|\) roots, so it is zero, and \(R_e(v')=0\) for every \(e\). By induction every
\(R_e\) is zero.
\end{proof}

Mathlib proves this lemma as part of the combinatorial Nullstellensatz, and
the Lean proof uses that version.

\begin{lemma}[Grid reduction]\label{lem:grid-reduction}
Let \(D\) be an integral domain, \(V\subseteq D\) finite,
\(L_V=\prod_{v\in V}(Y-v)\), and \(\rho_N\) the remainder of \(Y^N\) modulo \(L_V\).
Let \(I\) be finite, \(c_i\in D\) and \(d_{i,a}\ge0\) for \(i\in I\) and
\(a=1,\ldots,h\). If
\[
 \sum_{i\in I}c_i\prod_{a=1}^hv_a^{d_{i,a}}=0\qquad\text{for all }v\in V^h,
\]
then for all exponents \(e_1,\ldots,e_h<|V|\),
\[
 \sum_{i\in I}c_i\prod_{a=1}^h\bigl[Y^{e_a}\bigr]\rho_{d_{i,a}}=0
\]
\lean{claims/PerOrderSubspaceReduction.lean}.
\end{lemma}

\begin{proof}
If \(V\) is empty, the conclusion is vacuous for \(h\ge1\) and is the hypothesis
for \(h=0\); so let \(V\) be nonempty. Let \(R=\sum_ic_i\prod_a\rho_{d_{i,a}}(Y_a)\). It has degree less than \(|V|\) in
each variable, because \(\deg\rho_N<|V|\). Since \(L_V\) divides
\(Y^N-\rho_N\) and vanishes on \(V\), \(\rho_N(v)=v^N\) for \(v\in V\), so \(R\)
vanishes on \(V^h\). By Lemma~\ref{lem:grid-vanishing}, \(R=0\). Its coefficient of
\(\prod_aY_a^{e_a}\) is the displayed sum.
\end{proof}

\subsection{The lower bound}

\begin{theorem}\label{thm:per-order-lower}
Let \(n\ge1\) and \(0\le\ell<k\). Every \(P\in\Ft[x_\sigma]\) with
\(\mult_a(P)\ge k\) for all nonzero \(a\) and \(\mult_{\zero}(P)=\ell\) satisfies
\(\deg P\ge\Phi(n,k,\ell)\) \lean{claims/PerOrderLowerBound.lean}.
\end{theorem}

\begin{proof}
Let \(m=k-\ell-1=q2^n+r\) with \(0\le r<2^n\), and \(h=2q+s_2(r)\). By
Lemma~\ref{lem:legendre}, \(\Phi(n,k,\ell)=n+2(k-1)-h\), and
\(m=\sum_{a=1}^h2^{b_a}\) with \(b_a\le n-1\) for all \(a\). Since every part is at
least one, \(h\le m\le k-1\), and \(2^{b_a}\le m\) gives \(b_a<k\).

Suppose that \(\deg P\le\Phi(n,k,\ell)-1\). Take \(K=k\), so \(2^K\ge k\), an
expansion of \(P\) (Lemma~\ref{lem:expansion}), and \(w=P(\hat x)\) as in
Lemma~\ref{lem:congruence}. Form \(E^{(t)}_j\) from this \(w\).

\emph{Step 1: the conditions \((C_t)\).} For \(0\le t<k\) and \(T\subseteq\sigma\)
with \(|T|+2t\le h\),
\[
 \Bigl(w\prod_{i\in T}\hat x_i\Bigr)_{k-1-t}=0 .
\]
First let \(S\subseteq\sigma\) with \(|S|+2t\le h\), and \(\varepsilon=\sigma\setminus S\).
A monomial \(y^e\) with \(|e|=k-1-t\) in \(A_\varepsilon\) would have
\(|\varepsilon|+2|e|=n-|S|+2(k-1)-2t\ge n+2(k-1)-h>\deg P\), contradicting
Lemma~\ref{lem:expansion}. So \((A_\varepsilon)_{k-1-t}=0\), and
Lemma~\ref{lem:congruence} gives
\((w\prod_{i\in S}(1+\hat x_i))_{k-1-t}=0\). Expanding
\(\prod_{i\in S}(1+\hat x_i)=\sum_{T\subseteq S}\prod_{i\in T}\hat x_i\), and using that
every \(T\subseteq S\) also satisfies \(|T|+2t\le h\), induction on \(|S|\) gives the
claim for \(T=S\).

\emph{Step 2: vanishing on the span.} We show, by induction on \(j\) and for all
\(t\) at once, that \(E^{(t)}_j(y_{i_1},\ldots,y_{i_j})=0\) for all
\(i_1,\ldots,i_j\in\sigma\) whenever \(j+2t\le h\). Then \(t\le h/2<k\).
\begin{itemize}
 \item If the indices are distinct, this is Step~1 with \(T=\{i_1,\ldots,i_j\}\), by
 Lemma~\ref{lem:forms}(2). This covers \(j\le1\).
 \item Otherwise, by symmetry, the tuple is \((y_i,y_i,Z)\), and by
 Lemma~\ref{lem:forms}(3) the value is
 \(E^{(t)}_{j-1}(y_i,Z)+y_i\,E^{(t+1)}_{j-2}(Z)\). Both terms vanish by induction,
 since \((j-1)+2t\le h\) and \((j-2)+2(t+1)\le h\).
\end{itemize}
Let \(V=\{\sum_{i\in U}y_i:U\subseteq\sigma\}\), the \(\Ft\)-span of the variables, with
\(|V|=2^n\). By additivity (Lemma~\ref{lem:forms}(1)), \(E^{(0)}_h\) vanishes on
\(V^h\).

\emph{Step 3: reduction.} Apply Lemma~\ref{lem:grid-reduction} with
\(D=\Ft[y_\sigma]\), this \(V\), \(I=\{0,\ldots,k-1\}^h\),
\(c_\tau=w_{k-1-\sum_a2^{\tau_a}}\), \(d_{\tau,a}=2^{\tau_a}\), and
\(e_a=2^{b_a}<2^n\). It gives
\[
 \sum_\tau w_{k-1-\sum_a2^{\tau_a}}\prod_a\bigl[Y^{2^{b_a}}\bigr]\rho_{2^{\tau_a}}=0 .
\]
Since \(L_V\) is monic of degree \(2^n\), \(\rho_N=Y^N\) for \(N<2^n\).
\begin{itemize}
 \item If \(\tau_a<b_a\) for some \(a\), then \(\rho_{2^{\tau_a}}=Y^{2^{\tau_a}}\) has
 no \(Y^{2^{b_a}}\) term, and the term vanishes.
 \item If \(\tau_a\ge b_a\) for all \(a\) and \(\tau\ne b\), then
 \(\sum_a2^{\tau_a}>m\), so the index \(k-1-\sum_a2^{\tau_a}\) is below \(\ell\), and
 \(w\) has no such component by Lemma~\ref{lem:lowest-part}.
 \item For \(\tau=b\), \(\rho_{2^{b_a}}=Y^{2^{b_a}}\), and the term is
 \(w_{k-1-m}=w_\ell\).
\end{itemize}
So \(w_\ell=0\), contradicting Lemma~\ref{lem:lowest-part}.
\end{proof}

In the construction, each block of \(2^n\) in \(m\) is paid by a factor \(F^2\);
correspondingly, each block contributes two parts \(2^{n-1}\) to the
decomposition of \(m\) used in this proof.
